\documentclass[11pt]{article}
\usepackage{amsmath,amssymb,amsthm}
\usepackage{booktabs}
\usepackage{enumitem}
\usepackage{hyperref}
\usepackage[margin=1in]{geometry}

\hypersetup{colorlinks=true,linkcolor=black,citecolor=blue,urlcolor=blue}

\newtheorem{theorem}{Theorem}
\newtheorem{lemma}[theorem]{Lemma}
\newtheorem{corollary}[theorem]{Corollary}
\newtheorem{definition}{Definition}
\newtheorem{remark}{Remark}
\newtheorem{proposition}[theorem]{Proposition}

\title{Zoo Adopts the Liquidity Protocol\\
{\large Composition of Zoo NFT-LP (2021), Per-LLM Chains, and DAO
with Liquidity.io's Liquidity Protocol (2026)}}
\author{Zoo Labs Foundation\\
\texttt{research@zoo.ngo}}
\date{Adoption: 2026-04-20}

\begin{document}
\maketitle

\begin{abstract}
On 2026-04-20 Zoo Labs Foundation formally adopts the Liquidity
Protocol (Liquidity.io, 2026-04-01) as the canonical regulated
securities-collateralised settlement layer for treasury and
contributor compensation, while explicitly retaining the
2021 Zoo trade-name ``NFT Liquidity Protocol'' for
NFT-collateralised lending. We prove formal differentiation
(disjoint asset classes) and composition (both stacks compose under
Quasar 3.0 finality without interference). The result is that Zoo
operates two distinct ``liquidity protocols'' on the same chain
with no semantic collision and no cross-contamination of state.
\end{abstract}

\tableofcontents

%% ==================================================================
\section{Two protocols, one trade name}
\label{sec:two}
%% ==================================================================

\begin{table}[h]
\centering
\caption{Trade-name disambiguation}
\begin{tabular}{lll}
\toprule
& Zoo NFT-LP (2021) & Liquidity.io LP (2026) \\
\midrule
Asset class      & ERC-721 / ZRC-721         & ERC-3643 \\
Quote token      & ERC-20 stable             & USDL (1:1 USD) \\
Counterparty     & Zoo Lender Pool           & Alpaca Securities \\
Coined           & October 2021              & April 2026 \\
Citation         & \cite{zoo-2021}           & \cite{liquidity-lp} \\
Soundness proof  & \cite{zoo-nft-lp}         & \cite{liquidity-lp} \S5 \\
\bottomrule
\end{tabular}
\end{table}

\begin{remark}[Naming policy]
In Zoo prose: ``NFT Liquidity Protocol'' or ``Zoo NFT-LP'' for the
2021 trade name; ``Liquidity Protocol'' or ``Liquidity.io's LP'' for
the 2026 protocol. We never abbreviate either as just ``LP'' without
qualification.
\end{remark}

%% ==================================================================
\section{Adoption boundary}
\label{sec:adopt}
%% ==================================================================

\begin{definition}[Zoo adoption boundary]
\label{def:adopt}
``Zoo adopts the Liquidity Protocol'' means:
\begin{enumerate}[nosep]
  \item Zoo treasury disbursements (NGO payouts, DAO compensation)
        settle via the Liquidity Protocol's OMA on Lux C-Chain (or
        Zoo D-Chain when LP-134's D-Chain lands).
  \item Per-LLM-chain contributor compensation, when paid in
        securities-equivalent dollars, settles via the Liquidity
        Protocol.
  \item The 2021 Zoo NFT-LP continues to operate independently on
        Zoo L1 with its own state space and ABI; Zoo carries forward
        the 2021 trade-name claim under
        Proposition~\ref{prop:disjoint}.
\end{enumerate}
This is forward-only.  Pre-2026-04-20 NFT-LP operations are not
re-routed; Zoo treasury continues to support direct stablecoin
disbursement as a fallback.
\end{definition}

%% ==================================================================
\section{Disjointness}
\label{sec:disjoint}
%% ==================================================================

\begin{proposition}[Disjoint state spaces]
\label{prop:disjoint}
The state spaces are disjoint at every level:
\[
  \mathcal{S}_{\mathrm{Zoo\text{-}NFT\text{-}LP}} \cap
  \mathcal{S}_{\mathrm{Liquidity\text{-}LP}} = \emptyset.
\]
\begin{enumerate}[nosep]
  \item ABI selectors differ;
  \item contract addresses differ (different deployer accounts,
        different bytecode);
  \item asset classes differ
        (Prop.~\ref{prop:nft-vs-sec}).
\end{enumerate}
\end{proposition}

\begin{proposition}[NFT vs.\ security tokens]
\label{prop:nft-vs-sec}
ERC-721 (NFT) and ERC-3643 (security) are distinct EVM token
standards with distinct interfaces. No token is simultaneously
both. Even if a single token contract implemented both interfaces
(``hybrid token''), the dispatch is via function selector, so the
LP that handles it is determined by which interface is invoked.
\end{proposition}

\begin{proof}
Both follow from inspection of the standards. ERC-721 has unique
\texttt{tokenId}-keyed storage; ERC-3643 has compliance-gated
transfer. A single contract implementing both would route requests
through the appropriate selector; the runtime never confuses a
\texttt{transferFrom(uint256)} call (ERC-721) with a
\texttt{transfer(address,uint256)} call (ERC-3643).
\end{proof}

%% ==================================================================
\section{Composition theorem}
\label{sec:compose}
%% ==================================================================

\begin{theorem}[Zoo composition under Quasar finality]
\label{thm:compose}
Let $b$ be a Lux execution-chain block containing transactions
$T_b$ that interleave Zoo NFT-LP, Liquidity Protocol, Zoo per-LLM
chain commitments, and DAO governance ciphertexts. Under
\cite[Theorem 7.5]{quasar-cert} (QuasarCert soundness),
\cite[Theorem 5.2]{liquidity-lp} (LP atomicity), and
\cite[Theorem 4.1]{zoo-nft-lp} (NFT-LP soundness), the post-block
state $s_b'$ is:
\begin{enumerate}[nosep]
  \item deterministic given $s_{b-1}$ and $T_b$;
  \item satisfies all per-protocol invariants
        (NFT-LP atomic lending, LP atomic swap, DAO weighted-vote
        correctness);
  \item finalised under Quasar 3.0 PQ finality.
\end{enumerate}
\end{theorem}

\begin{proof}
By Prop.~\ref{prop:disjoint} the state updates of distinct
protocols are non-overlapping. Block-STM serialisability
(\cite{dagevm-lean}) ensures determinism. Each protocol's atomicity
property is preserved within a single transaction (per-protocol
proofs cited above). QuasarCert finalises the unique post-state by
\cite[Theorem 7.5]{quasar-cert}.
\end{proof}

\begin{corollary}[No cross-contamination]
\label{cor:no-cross}
A bug or exploit in Liquidity.io's LP cannot affect Zoo NFT-LP
state; a bug in Zoo NFT-LP cannot affect Liquidity.io LP state.
The two are isolated by Prop.~\ref{prop:disjoint}.
\end{corollary}

%% ==================================================================
\section{Detailed proof of Theorem \ref{thm:compose}}
\label{sec:proof-detail}
%% ==================================================================

We give a more careful argument.

\paragraph{Setup.}
Block $b$ contains a sequence of transactions
$T_b = (\mathsf{tx}_1, \ldots, \mathsf{tx}_N)$, each labelled with
its target protocol (NFT-LP, LP, per-LLM-chain, DAO, or
unaffiliated EVM). The block is built by some honest proposer; the
post-state is computed by Block-STM scheduling.

\paragraph{Step 1: per-protocol disjointness.}
By Prop.~\ref{prop:disjoint}, the read/write sets of distinct
protocols do not intersect at the storage-key level (different
contract addresses imply different storage prefixes). Block-STM's
optimistic scheduler thus never re-executes a transaction due to
cross-protocol conflict.

\paragraph{Step 2: per-protocol atomicity.}
Each transaction is single-EVM-tx atomic: failures inside the call
graph revert the whole transaction. Hence within each protocol,
the per-protocol invariants (NFT-LP atomic lending,
LP atomic swap, DAO single-step proposal closure, per-LLM-chain
single-step gradient apply) hold.

\paragraph{Step 3: serialisability.}
Block-STM produces a serial schedule $\pi$ (\cite{dagevm-lean}).
Running $\pi$ on $s_{b-1}$ yields $s_b'$. By determinism of EVM
and pinned FP modes (per-LLM-chain only), $s_b'$ is unique.

\paragraph{Step 4: finality.}
QuasarCert $\mathcal{C}_b$ binds to
$\mathsf{subj}_b$ which commits to the post-state root $s_b'$. By
\cite[Theorem 7.5]{quasar-cert} the cert is sound: an adversary
cannot finalise a different post-state.

\paragraph{Result.}
$s_b'$ is deterministic, satisfies all per-protocol invariants,
and is PQ-finalised. \qed

%% ==================================================================
\section{Worked example: NGO disbursement}
\label{sec:example}
%% ==================================================================

A passing DAO proposal disburses 10{,}000 USDL to WWF.  The flow:

\begin{enumerate}[nosep]
  \item DAO contract emits intent
        $(t = \mathrm{now}, P = \mathrm{proposal}\#42,
          \mathrm{NGO}=\text{WWF}, \mathrm{amount}=10{,}000)$
        after Theorem~5.1 of \cite{zoo-dao} (tally correctness)
        accepts.
  \item Treasury bot constructs OMA call:
        $\texttt{OMA.swap(USDL, USDL, 10000, recipient=WWF\_addr)}$
        on Lux C-Chain.
  \item The OMA, observing same-token swap, treats this as a
        direct transfer; ATS inventory check passes; transfer
        executes.
  \item Block finalises under Quasar 3.0; WWF wallet receives
        $10{,}000$ USDL minted from ATS-backed reserves.
\end{enumerate}

By Theorem~\ref{thm:compose} this commits atomically with the DAO
state (proposal marked disbursed) and is PQ-final.

%% ==================================================================
\section{Worked example: NFT-LP and Liquidity-LP coexistence}
\label{sec:coexist}
%% ==================================================================

A Lux C-Chain block $b$ contains:
\begin{itemize}[nosep]
  \item Tx 1: NFT-LP loan issuance --- borrower deposits a Zoo
        animal NFT, receives $1{,}000$ ZUSD;
  \item Tx 2: NFT-LP marketplace trade --- buyer purchases an NFT
        for $500$ ZUSD; $12.50$ goes to sustainability treasury
        per Lemma~\ref{lem:tax} of \cite{zoo-nft-lp};
  \item Tx 3: Liquidity-Protocol OMA swap --- user converts
        $50$ USDL to $0.21$ AAPL (via ATS inventory).
\end{itemize}

After Block-STM execution, all three transactions commit. By
Proposition~\ref{prop:disjoint}, no cross-contamination is possible:
the NFT-LP state map and Liquidity-LP state map have disjoint keys
(ABI selectors and contract addresses differ). By
Theorem~\ref{thm:compose}, the post-state is unique and
PQ-finalised.

%% ==================================================================
\section{Operational notes}
\label{sec:ops}
%% ==================================================================

\paragraph{Treasury bot.}
The Zoo treasury bot is a permissionless on-chain script. It
reads DAO outcomes and emits OMA settlement transactions; it holds
no privileged keys (the DAO multisig is the only authority).

\paragraph{Naming in code.}
Solidity / Move / Vyper code in the Zoo repos uses \texttt{NFTLP}
(no spaces) for the 2021 trade name and
\texttt{LiquidityProtocol} (camel-case) for the 2026 Liquidity.io
adoption. The two are referenced via separate import paths to
prevent any namespace collision.

\paragraph{Brand mapping.}
Both protocols operate within the Zoo brand surface; the Zoo
documentation site distinguishes them by visual styling
(NFT-LP: animal-themed; LP: Liquidity.io blue). No protocol-level
ambiguity arises.

%% ==================================================================
\section{Status, 2026-04-20}
\label{sec:status}
%% ==================================================================

\begin{itemize}[nosep]
  \item Zoo L1: live (with Zoo D-/E-/F-Chain white-label resells of
        Lux templates).
  \item Zoo NFT-LP: live (V4 rebuild, Lux C-Chain).
  \item Zoo per-LLM chains: live (Block-STM-as-consensus on Lux).
  \item Zoo DAO: live (homomorphic tally, M-Chain MPC threshold
        decrypter).
  \item Liquidity.io LP: live since 2026-04-01.
  \item Cross-link: Zoo treasury disbursement contract upgraded
        2026-04-20 to support OMA settlement (default-on for new
        proposals).
\end{itemize}

%% ==================================================================
\begin{thebibliography}{9}

\bibitem{zoo-2021}
A.~Worring, Z.~Kelling.
\newblock Zoo: NFT-Driven Conservation, October 2021.
\newblock \texttt{zooai/papers/zoo-2021-original-whitepaper/}.

\bibitem{liquidity-lp}
Liquidity.io.
\newblock The Liquidity Protocol v1.0.
\newblock \texttt{liquidityio/proofs/liquidity-protocol/}, April 2026.

\bibitem{zoo-nft-lp}
Zoo Labs Foundation.
\newblock Zoo NFT Liquidity Protocol: Soundness.
\newblock \texttt{zooai/proofs/zoo-nft-liquidity-protocol-soundness.tex}, 2026.

\bibitem{quasar-cert}
Z.~Kelling.
\newblock QuasarCert Soundness, Quasar 3.0.
\newblock \texttt{luxfi/proofs/quasar-cert-soundness.tex}, 2025-12-15.

\bibitem{dagevm-lean}
Lux Industries.
\newblock \texttt{lean/Consensus/DAGEVM.lean}.

\bibitem{lux-adopts}
Lux Industries.
\newblock Lux Adopts the Liquidity Protocol.
\newblock \texttt{luxfi/proofs/lux-adopts-liquidity-protocol.tex}, 2026-04-20.

\end{thebibliography}

\end{document}
